\documentclass{article}
\usepackage[utf8]{inputenc}
\usepackage[T1]{fontenc}
\usepackage{geometry}
\usepackage{amsmath}
\usepackage{listings}
\usepackage{booktabs}
\usepackage{hyperref}
\usepackage{xcolor}

\geometry{a4paper, margin=1in}

\lstset{
    language=Python,
    basicstyle=\small\ttfamily,
    breaklines=true,
    frame=single,
    showstringspaces=false,
    keywordstyle=\color{blue},
    commentstyle=\color{gray},
    stringstyle=\color{orange}
}

\title{Graph Diversification: Potensi Pengembangan ke Machine Learning}
\author{Catatan Tesis}
\date{\today}

\begin{document}

\maketitle

\section*{Graph Diversification: Bisa Dijadikan Machine Learning?}
Jawaban singkat: \textbf{Ya, bisa} — dan ini bisa jadi kontribusi tesis Anda yang paling kuat!

\section*{Memahami Graph Diversification Saat Ini}
Cara kerjanya sekarang (bukan ML, ini \textit{rule-based}):

\begin{center}
    Data Return Historis $\rightarrow$ Hitung Matriks Korelasi (matematika biasa) $\rightarrow$ Bangun Graph (edge = korelasi $> \theta$) $\rightarrow$ Cari Maximum Independent Set / MIS (algoritma deterministik) $\rightarrow$ Beri bobot sama rata ke aset terpilih.
\end{center}

\textbf{Masalahnya}: Tidak ada yang ``dipelajari'' — threshold $\theta$ ditentukan manual (0.4 atau 0.5), bobot selalu sama rata, tidak ada adaptasi dari data.

\section*{Evolusi ke Machine Learning: 3 Level}

\subsection*{Level 1 — Belajar Threshold (Sederhana)}
\begin{itemize}
    \item \textbf{Saat ini}: $\theta = 0.4$ (manual, \textit{fixed}).
    \item \textbf{Versi ML}: $\theta^* = \arg\max \text{Sharpe}(\theta)$ melalui \textit{cross-validation} atau \textit{grid search}.
    \item Ini sudah bisa disebut ``machine learning'' karena parameter dipelajari dari data.
\end{itemize}

\subsection*{Level 2 — Belajar Bobot Aset (Menengah)}
\begin{itemize}
    \item \textbf{Saat ini}: Bobot = $1/n$ (sama rata untuk semua aset terpilih).
    \item \textbf{Versi ML}: Bobot = $f(\text{embedding GNN})$. Model memprediksi bobot optimal.
    \item Setiap node (aset) dalam graph menghasilkan \textit{embedding} dari GNN. \textit{Embedding} ini dipakai untuk prediksi bobot, bukan sekadar $1/n$.
\end{itemize}

\subsection*{Level 3 — End-to-End Learnable Portfolio (Canggih)}
\begin{itemize}
    \item \textbf{Input}: Return historis + Struktur Graph.
    \item \textbf{GNN Layer}: Belajar representasi aset.
    \item \textbf{Attention Layer}: Belajar MANA aset yang paling ``independen''.
    \item \textbf{Softmax $\rightarrow$ Bobot Portofolio}.
    \item \textbf{Loss Function}: $-\text{Sharpe Ratio}$ (\textit{maximize}).
    \item \textbf{Backpropagation}: \textit{Update} semua parameter. Ini adalah Graph Portfolio Network yang sepenuhnya \textit{end-to-end}.
\end{itemize}

\section*{Perbandingan Pendekatan}

\begin{table}[h]
    \centering
    \begin{tabular}{lll}
        \toprule
        \textbf{Pendekatan} & \textbf{Apa yang Dipelajari} & \textbf{Nama Teknis} \\
        \midrule
        GD klasik (sekarang) & Tidak ada & Rule-based \\
        GD + threshold learning & Nilai $\theta$ optimal & Grid Search / BO \\
        GD + GNN weights & Bobot dari embedding & GNN Regression \\
        GD + Attention & Seleksi aset + bobot & Graph Attention Network (GAT) \\
        GD + RL & Keputusan rebalancing & Deep RL / DQN \\
        End-to-end & Segalanya & E2E Portfolio GNN \\
        \bottomrule
    \end{tabular}
\end{table}

\section*{Yang Paling Realistis untuk Tesis Anda}
Dari hasil eksperimen, GD ($\theta=0.5$) adalah strategi terbaik (Sharpe 1.09 di \textit{recovery}!). Langkah natural berikutnya adalah menggantikan bobot sama rata dengan bobot yang dipelajari GNN:

\begin{lstlisting}[caption=Perbandingan Kode GD Klasik vs ML]
# Saat ini (GD klasik):
weights[asset] = 1.0 / n_selected   # <- sama rata

# Versi ML yang diusulkan:
embedding = gnn(features, edges)     # <- GNN belajar representasi
weights = softmax(linear(embedding)) # <- belajar bobot optimal
# Dilatih untuk maximize Sharpe Ratio
\end{lstlisting}

\section*{Keunggulan Pendekatan Ini untuk Tesis}
\begin{itemize}
    \item[\checkmark] Berbasis pada temuan nyata (GD sudah terbukti bagus).
    \item[\checkmark] Menambah komponen ML yang \textit{genuine} (bukan sekedar estimasi korelasi).
    \item[\checkmark] Dapat dibandingkan langsung dengan GD klasik sebagai \textit{baseline}.
    \item[\checkmark] Kontribusi jelas: ``\textit{ML-enhanced Graph Diversification vs. Rule-based Graph Diversification}''.
\end{itemize}

\section*{Apakah Perlu Diimplementasikan Sekarang?}
Tidak harus. Untuk tesis, Anda bisa:
\begin{enumerate}
    \item \textbf{Implementasi penuh}: Butuh waktu lebih, tapi kontribusi lebih besar.
    \item \textbf{Sebagai ``Future Work''}: Tulis di bab kesimpulan bahwa ini arah penelitian selanjutnya yang menjanjikan.
\end{enumerate}

Mengingat GD ($\theta=0.5$) sudah sangat bagus (Sharpe 1.09), argumen tesis yang kuat adalah: ``Graph Diversification berbasis \textit{rule} terbukti unggul. Pengembangan lebih lanjut dengan ML (\textit{learned weights}) berpotensi meningkatkan performa lebih jauh.''

\end{document}
